
\chapter{Herramientas de seguridad implementadas}

Las herramientas se eligieron a partir de las amenazas del modelo STRIDE:
cada una debía cubrir al menos una amenaza, funcionar con el stack del
proyecto (Node.js, Express, Vite, Docker y PostgreSQL), ejecutarse en GitHub
Actions o en local sin infraestructura adicional y ser gratuita. Los scripts
de ejecución se encuentran en \texttt{scripts/security/} y generan sus
reportes en \texttt{reports/}. Los resultados que se presentan en este
capítulo provienen de la ejecución real del pipeline en GitHub Actions
(23 de noviembre de 2026, rama
\url{feature/devsecops-threat-model-sast-dast}) y de los reportes de ZAP
generados en local.

% ============================================================
% SAST
% ============================================================

\section{Análisis estático de seguridad (SAST)}

\subsection{Herramienta seleccionada}

Se utilizó \textbf{Semgrep} en su versión de línea de comandos con reglas de
la comunidad. Se eligió porque es rápido, tiene reglas para JavaScript,
TypeScript, Dockerfile y archivos de configuración, no requiere cuenta ni
servidor propio y se ejecuta en cada \textit{push}. Su objetivo es detectar
patrones inseguros en el código fuente antes de que lleguen a ejecución.

En la fase de selección se evaluó también SonarCloud como complemento, pero
no se integró en el pipeline actual; Semgrep es el único análisis estático
implementado.

\subsection{Configuración e integración}

El script \texttt{scripts/security/run-semgrep.sh} ejecuta la imagen oficial
\texttt{semgrep/semgrep} en un contenedor, con tres conjuntos de reglas:
\texttt{p/ci}, \texttt{p/security-audit} y \texttt{p/javascript}. Se
desactivan las métricas (\texttt{-{}-metrics=off}), se excluye la carpeta
\texttt{reports} y el resultado se guarda en formato JSON en
\texttt{reports/sast/semgrep-report.json}.

En GitHub Actions, el job \texttt{semgrep-scan} ejecuta el script después de
los jobs de detección de secretos (\texttt{needs}) y publica el reporte como
\textit{artifact}. El paso corre con \texttt{continue-on-error}, de modo que
un hallazgo no impide que el reporte llegue a DefectDojo.

\begin{lstlisting}[language=yaml, caption={Job de Semgrep en el pipeline.}, label={lst:semgrep}]
semgrep-scan:
  needs: [gitleaks-scan, trufflehog-scan]
  runs-on: ubuntu-latest
  steps:
    - uses: actions/checkout@v4
    - name: Semgrep
      continue-on-error: true
      run: bash scripts/security/run-semgrep.sh
    - uses: actions/upload-artifact@v4
      if: always()
      with:
        name: reports-semgrep
        path: reports/sast/semgrep-report.json
\end{lstlisting}

\subsection{Resultados obtenidos}

Semgrep analizó 68 archivos versionados con 122 reglas y reportó
\textbf{15 hallazgos}, todos con la misma regla
(\texttt{github-actions-mutable-action-tag}, severidad \textit{WARNING},
CWE-1357). Los 15 corresponden a acciones de GitHub referenciadas por una
etiqueta mutable (por ejemplo \texttt{actions/checkout@v4}) en el archivo
\texttt{.github/workflows/pipeline.yml}, en lugar de fijarlas por el
\textit{hash} del commit. Es un riesgo de cadena de suministro: si la
etiqueta se redirige a código malicioso, el pipeline lo ejecutaría.

\textbf{No se detectaron hallazgos en el código de la aplicación}
(\texttt{backend/src} y \texttt{frontend/src}). Esto no significa que el
código esté libre de problemas: las reglas estándar no detectaron el CORS
abierto (\texttt{cors()}) ni la devolución de \texttt{err.message}, que sí
están confirmados por otras vías (capítulo de modelado de amenazas). Para
cubrirlos haría falta escribir reglas propias de Semgrep, lo cual queda como
mejora pendiente.

% ============================================================
% DAST
% ============================================================

\section{Análisis dinámico de seguridad (DAST)}

\subsection{Herramienta seleccionada}

Se utilizó \textbf{OWASP ZAP} en modo \textit{baseline}. ZAP es gratuito, de
código abierto, dispone de imágenes oficiales para Docker y su escaneo
\textit{baseline} es \textbf{pasivo}: observa las respuestas de la aplicación
sin enviar ataques, por lo que es seguro ejecutarlo contra el entorno local.

\subsection{Configuración e integración}

El script \texttt{scripts/security/run-zap-baseline.sh} ejecuta la imagen
\texttt{zaproxy/zap-stable} dentro de la misma red de Docker Compose que la
aplicación, por lo que el objetivo se indica con el nombre del servicio
(\texttt{frontend}). Se realizaron dos escaneos: uno sobre la aplicación web
y otro sobre el endpoint de la API (\texttt{/api/notes}). Cada uno genera
reportes en HTML, JSON y XML en \texttt{reports/dast/}.

{\sloppy
Las alertas se evalúan con el archivo \texttt{zap/rules.conf}, que define
una puerta de calidad (\textit{gate}): las alertas de \textit{Content
Security Policy} ausente, \textit{Anti-clickjacking}, \textit{X-Content-Type-Options},
\textit{Strict-Transport-Security} y \textit{Cross-Domain Misconfiguration}
tienen acción \texttt{FAIL}; otras, como \textit{Permissions Policy}, solo
generan aviso (\texttt{WARN}); y el contenido cacheable se ignora.
\par}

Este análisis \textbf{se ejecuta en local} y requiere que la aplicación esté
levantada con \texttt{docker compose}. Todavía no forma parte del pipeline de
GitHub Actions, porque exigiría levantar el stack completo dentro del
\textit{runner}.

\subsection{Resultados obtenidos}

\begin{table}[H]
\centering
\caption{Alertas de ZAP por nivel de riesgo (tipos de alerta).}
\label{tab:zap-resumen}

\begin{tabularx}{\textwidth}{l c c c c c}
\toprule
\textbf{Escaneo} &
\textbf{Alto} &
\textbf{Medio} &
\textbf{Bajo} &
\textbf{Informativo} &
\textbf{Total} \\
\midrule
Aplicación web & 1 & 2 & 7 & 3 & 13 \\
API (\texttt{/api/notes}) & 1 & 4 & 9 & 3 & 17 \\
\bottomrule
\end{tabularx}
\end{table}

Las alertas más relevantes y su relación con el modelo de amenazas son:

\begin{table}[H]
\centering
\caption{Alertas principales de ZAP y amenaza asociada.}
\label{tab:zap-alertas}

\begin{tabularx}{\textwidth}{l X l}
\toprule
\textbf{Riesgo} &
\textbf{Alerta} &
\textbf{Amenaza} \\
\midrule
Alto & \textit{nginx Version with Known Vulnerabilities}: el encabezado
\texttt{Server} revela una versión de Nginx con vulnerabilidades conocidas. &
T4 \\

Medio & \textit{Content Security Policy (CSP) Header Not Set}. & T4 \\

Medio & \textit{Missing Anti-clickjacking Header}. & T4 \\

Medio & \textit{Cross-Domain Misconfiguration}: la API responde
\texttt{Access-Control-Allow-Origin: *}. & T3 \\

Bajo & \textit{Application Error Disclosure}: respuesta HTTP 500 en
\texttt{/api/notes}. & I1 \\

Bajo & \textit{X-Powered-By} y \textit{Server} revelan Express y la versión de
Nginx. & I1 \\

Bajo & \textit{X-Content-Type-Options Header Missing}. & T4 \\
\bottomrule
\end{tabularx}
\end{table}

Las alertas con acción \texttt{FAIL} (CSP, anti-clickjacking,
\textit{X-Content-Type-Options} y \textit{Cross-Domain Misconfiguration})
aparecen en los resultados, por lo que el \textit{gate} de ZAP no se
supera. ZAP confirma así tres amenazas del modelo (T3, T4 e I1) y añade dos
hallazgos que el modelo no había identificado: la versión de Nginx vulnerable
y la exposición de la tecnología del servidor.

La figura~\ref{fig:dast} muestra el resumen del reporte de ZAP para el
escaneo de la API.

\evidencefigure
{figures/dast-reporte.png}
{Resumen del reporte de OWASP ZAP para la API: 1 alerta alta, 4 medias, 9 bajas y 3 informativas.}
{fig:dast}

% ============================================================
% SECRET SCANNING
% ============================================================

\section{Detección de secretos (Secret Scanning)}

\subsection{Herramienta seleccionada}

Se utilizaron dos herramientas complementarias:

\begin{itemize}
    \item \textbf{Gitleaks:} escanea el historial completo de Git en busca
    de credenciales, claves de API y tokens mediante reglas y expresiones
    regulares. Es rápida y no necesita cuenta.
    \item \textbf{TruffleHog:} escanea los archivos del sistema de archivos
    y, cuando es posible, verifica contra el proveedor si una credencial
    sigue activa. Se configuró para reportar resultados
    \texttt{verified} y \texttt{unknown}.
\end{itemize}

El secreto más sensible del proyecto es la contraseña de PostgreSQL
(\texttt{pg\_password.txt}). Ambas herramientas permiten comprobar que no
se haya versionado por error.

\subsection{Configuración e integración}

{\sloppy
Ambos escaneos se ejecutan en contenedores Docker mediante
\texttt{run-gitleaks.sh} y \texttt{run-truffleHog.sh}. Gitleaks usa el
comando \texttt{detect} sobre todo el repositorio y TruffleHog el comando
\texttt{filesystem}, excluyendo \texttt{.git}, \texttt{reports} y
\texttt{node\_modules} mediante \texttt{trufflehog-exclude.txt}. En GitHub
Actions, los jobs \texttt{gitleaks-scan} y \texttt{trufflehog-scan} se
ejecutan con el historial completo (\texttt{fetch-depth: 0}) y publican sus
reportes como \textit{artifacts}.
\par}

\subsection{Resultados obtenidos}

\textbf{Ninguna de las dos herramientas encontró secretos.} Gitleaks informó
``no leaks found'' y devolvió un reporte JSON vacío; el reporte de
TruffleHog también quedó vacío. Esto es coherente con los controles del
proyecto: \texttt{pg\_password.txt} está en \texttt{.gitignore} y la
contraseña se entrega como \textit{Docker secret}, fuera de las imágenes. Por
ello no hay valores de credenciales que ocultar en esta sección.

% ============================================================
% SCA
% ============================================================

\section{Análisis de composición de software (SCA)}

\subsection{Herramienta seleccionada}

Se utilizó \textbf{Trivy}, que analiza dependencias a partir de los archivos
de bloqueo (\textit{lockfiles}) y compara cada versión con bases de datos de
vulnerabilidades conocidas. Se eligió por no requerir cuenta ni enviar los
manifiestos a un servicio externo, y porque con el mismo binario puede
ampliarse a imágenes de contenedor y configuración.

\subsection{Configuración e integración}

El script \texttt{scripts/security/run-trivy-fs.sh} ejecuta \texttt{trivy fs}
sobre el repositorio, con el escáner de vulnerabilidades
(\texttt{-{}-scanners vuln}) y filtrando por severidad alta y crítica. Analiza
los \texttt{package-lock.json} de \texttt{backend} y \texttt{frontend}. El job
\texttt{trivy-scan} del pipeline verifica antes que exista el archivo de
bloqueo, porque sin él Trivy no detecta las dependencias, y publica el
reporte JSON como \textit{artifact}.

\subsection{Resultados obtenidos}

Trivy encontró \textbf{una vulnerabilidad de severidad alta} en las
dependencias del backend:

\begin{table}[H]
\centering
\caption{Vulnerabilidad detectada por Trivy.}
\label{tab:trivy}

\begin{tabularx}{\textwidth}{l l l l l X}
\toprule
\textbf{CVE} &
\textbf{Paquete} &
\textbf{Instalada} &
\textbf{Corregida en} &
\textbf{Severidad} &
\textbf{Ubicación} \\
\midrule
CVE-2026-111181 & \texttt{body-parser} & 1.20.8 & 1.20.9 y 2.4.0 &
Alta & Backend (lockfile) \\
\bottomrule
\end{tabularx}
\end{table}

El título del aviso es \textit{Denial of Service via unhandled decompression
stream error}. \texttt{body-parser} es una dependencia de Express que
procesa el cuerpo de las peticiones, por lo que la vulnerabilidad afecta
directamente a la API. El \texttt{package-lock.json} del frontend no presentó
vulnerabilidades altas ni críticas. El problema se corrige actualizando
\texttt{body-parser} a la versión 1.20.9.

% ============================================================
% COMPARACIÓN
% ============================================================

\section{Resumen comparativo de herramientas}

\begin{table}[H]
\centering
\caption{Resumen de herramientas de seguridad.}
\label{tab:herramientas}

\begin{tabularx}{\textwidth}{l l X X}
\toprule
\textbf{Categoría} &
\textbf{Herramienta} &
\textbf{Objetivo} &
\textbf{Integración} \\
\midrule
SAST & Semgrep &
Patrones inseguros en código y configuración. 15 hallazgos, todos en el
pipeline. &
Job \texttt{semgrep-scan} en GitHub Actions \\

DAST & OWASP ZAP (baseline) &
Cabeceras, CORS y errores en la aplicación en ejecución. 13 alertas en la
web y 17 en la API. &
Local, con \texttt{run-zap-baseline.sh}; aún no está en el pipeline \\

Secret Scanning & Gitleaks y TruffleHog &
Credenciales en el historial y en los archivos. Sin hallazgos. &
Jobs \texttt{gitleaks-scan} y \texttt{trufflehog-scan} \\

SCA & Trivy &
Vulnerabilidades en dependencias npm. 1 hallazgo alto (\texttt{body-parser}). &
Job \texttt{trivy-scan} \\
\bottomrule
\end{tabularx}
\end{table}
